\documentclass[UTF8,11pt]{ctexart}
\usepackage[a4paper,margin=1.8cm]{geometry}
\usepackage{amsmath,amssymb,booktabs,longtable,array}
\usepackage{xcolor,hyperref,enumitem,fancyhdr,listings}
\hypersetup{colorlinks=true,linkcolor=blue!45!black,urlcolor=blue!45!black}
\setlength{\parindent}{0pt}
\setlength{\parskip}{0.35em}
\linespread{1.08}
\setcounter{secnumdepth}{0}
\setlist[itemize]{leftmargin=2em,itemsep=0.15em,topsep=0.2em}
\setlist[enumerate]{leftmargin=2.2em,itemsep=0.15em,topsep=0.2em}
\pagestyle{fancy}
\setlength{\headheight}{14pt}
\fancyhf{}
\lhead{CS338 Review Note}
\rhead{\thepage}
\lstset{basicstyle=\ttfamily\small,breaklines=true,columns=fullflexible,frame=single}
\title{CS338 Introduction to Theory of Computation 详细复习 Note}
\author{}
\date{\today}
\begin{document}
\maketitle
\tableofcontents
\newpage


资料覆盖：\texttt{ppts/} 中全部课件，\texttt{Week1} 至 \texttt{Week13} 的作业题面与 LaTeX/PDF 答案，\texttt{answers/} 中 Assignment 1-13 官方答案。  

使用方式：先看第 0 节答题规范，再按 DFA/NFA/RE -> CFL/PDA -> TM/Decidability -> Complexity 的顺序复习。


\bigskip\hrule\bigskip


\section{0. 考试答题规范：不要只写结论}


\subsection{0.1 Construction proof 必须写完整构造}


如果题目说 “use the construction proof”，答案必须给出新机器/新文法的完整形式，不能只画直觉图或说“显然封闭”。


\textbf{DFA 乘积构造模板（交、并、差、对称差）}


给定

\[
M_1=(Q_1,\Sigma,\delta_1,q_1,F_1),\quad
M_2=(Q_2,\Sigma,\delta_2,q_2,F_2).
\]

构造

\[
M=(Q_1\times Q_2,\Sigma,\delta,(q_1,q_2),F)
\]

其中

\[
\delta((r,s),a)=(\delta_1(r,a),\delta_2(s,a)).
\]

按目标语言设置接受态：


\begin{itemize}
\item 交集：\texttt{F=F\_1\textbackslash{}times F\_2}
\item 并集：\texttt{F=(F\_1\textbackslash{}times Q\_2)\textbackslash{}cup(Q\_1\textbackslash{}times F\_2)}
\item 差集：\texttt{F=F\_1\textbackslash{}times(Q\_2-F\_2)}
\item 对称差：\texttt{F=(F\_1\textbackslash{}times(Q\_2-F\_2))\textbackslash{}cup((Q\_1-F\_1)\textbackslash{}times F\_2)}
\end{itemize}


最后写一句正确性：读入同一前缀后，第一分量正是 \(M_1\) 的状态，第二分量正是 \(M_2\) 的状态，因此接受条件正好等价于目标布尔组合。


\textbf{NFA 并/连接/星号构造模板}


\begin{itemize}
\item 并：新起始态 \(q_0\)，从 \(q_0\) 用 \(\epsilon\)-transition 到两个旧起始态；接受态为两个 NFA 的接受态并集。
\item 连接：从第一个 NFA 的每个接受态用 \(\epsilon\)-transition 到第二个 NFA 的起始态；接受态为第二个 NFA 的接受态。
\item Star：新起始态也是接受态；从新起始态 \(\epsilon\) 到旧起始态；从每个旧接受态 \(\epsilon\) 回旧起始态。
\item 单接受态：新增唯一接受态 \(q_a\)，从所有旧接受态 \(\epsilon\) 到 \(q_a\)，原接受态取消接受。
\end{itemize}


\textbf{子集构造模板}


给 NFA \(N=(Q,\Sigma,\delta,q_0,F)\)，构造 DFA \(D=(2^Q,\Sigma,\delta',E(q_0),F')\)。  

若有 \(\epsilon\)-transition，则

\[
\delta'(S,a)=E\left(\bigcup_{q\in S}\delta(q,a)\right).
\]

接受态：

\[
F'=\{S\subseteq Q\mid S\cap F\neq\emptyset\}.
\]

作业里需要列出可达子集和转移表，不要只写“用 subset construction”。


\textbf{CFG -> PDA 模板}


构造 PDA 用栈模拟最左推导：


\begin{enumerate}
\item 初始先压入底符号 \texttt{\$}，再压入起始变量 \(S\)。
\item 对每条产生式 \(A\to u\)，加 \(\epsilon,A\to u^R\) 的替换路径（多符号压栈要拆成多个中间状态）。
\item 对每个终结符 \(a\)，加读入并匹配的转移 \(a,a\to\epsilon\)。
\item 栈底 \texttt{\$} 被弹出且输入读完时接受。
\end{enumerate}


\textbf{PDA -> CFG 模板}


先保证 PDA 只有一个接受态且接受前清空栈。变量 \(A_{pq}\) 表示从状态 \(p\) 到 \(q\) 并清空所压栈内容的所有字符串：


\begin{itemize}
\item 基础规则：\(A_{pp}\to\epsilon\)
\item 连接规则：\(A_{pq}\to A_{pr}A_{rq}\)
\item 匹配规则：若 \(p\xrightarrow{a,\epsilon\to t}r\) push \(t\)，且 \(s\xrightarrow{b,t\to\epsilon}q\) pop \(t\)，则
\[
  A_{pq}\to aA_{rs}b.
  \]

\end{itemize}


\textbf{归约证明模板}


要证明目标 \(B\) 难，必须从已知难问题 \(A\) 归约到 \(B\)：\(A\le_m B\) 或 \(A\le_p B\)。写法固定：


\begin{enumerate}
\item 说明已知 \(A\) 的性质，例如 \(A_{TM}\) undecidable，或 3SAT NP-complete。
\item 给出输入 \(x\) 到实例 \(f(x)\) 的构造，且构造可计算/多项式时间。
\item 证明双向：
\[
   x\in A \iff f(x)\in B.
   \]

\item 结论：若 \(B\) 可解，则 \(A\) 可解，矛盾；或 \(B\) NP-hard。若还要 NP-complete，必须再证明 \(B\in NP\)。
\end{enumerate}


\textbf{Pumping lemma 答题模板}


证明非正则：


\begin{enumerate}
\item 假设 \(L\) 正则，令 \(p\) 为 pumping length。
\item 选 \(w\in L\) 且 \(|w|\ge p\)。
\item 任取分解 \(w=xyz\)，满足 \(|xy|\le p, |y|>0\)。
\item 由 \(|xy|\le p\) 定位 \(y\) 在某一块内。
\item 选 \(i=0\) 或 \(i=2\)，证明 \(xy^iz\notin L\)，矛盾。
\end{enumerate}


证明非 CFL：


\begin{enumerate}
\item 假设 \(L\) 是 CFL，令 \(p\) 为 CFL pumping length。
\item 选 \(s\in L\)，分解 \(s=uvxyz\)，满足 \(|vxy|\le p, |vy|\ge 1\)。
\item 由于 \(|vxy|\le p\)，它最多跨越有限相邻块。
\item 分情况或统一说明 pumping 改变局部块，破坏全局相等/比例关系。
\item 选 \(i=0\) 或 \(i=2\)，得到 \(uv^ixy^iz\notin L\)。
\end{enumerate}


\bigskip\hrule\bigskip


\section{1. DFA 与正则语言}


\subsection{1.1 DFA 定义与接受}


DFA 是五元组

\[
M=(Q,\Sigma,\delta,q_0,F)
\]

其中 \(Q\) 是有限状态集，\(\Sigma\) 是字母表，\(\delta:Q\times\Sigma\to Q\)，\(q_0\) 是初始态，\(F\subseteq Q\) 是接受态。


DFA 接受 \(w=w_1\cdots w_n\)，当且仅当存在状态序列

\[
r_0=q_0,\quad r_{i+1}=\delta(r_i,w_{i+1}),\quad r_n\in F.
\]

语言 \(L(M)=\{w\mid M\text{ accepts }w\}\)。若某语言被某 DFA 识别，则它是 regular language。


\subsection{1.2 设计 DFA 的常用状态含义}


复习时不要背图，背“状态记录什么信息”：


\begin{itemize}
\item 以某串结尾：状态记录最长后缀匹配了目标串多少位。例如 contains \texttt{0101}，状态记录已匹配 \texttt{}, \texttt{0}, \texttt{01}, \texttt{010}, \texttt{0101}。
\item 至少/至多出现次数：状态记录计数并在阈值处吸收，例如至少四个 \texttt{1} 用 \(0,1,2,3,\ge4\) 五类。
\item 奇偶性：两个状态来回切换。
\item “以 0 开头且以 1 结尾”：先记住开头是否合法，再记录最后一个字符。
\item “不含某 substring”：先构造“含该 substring”的 DFA，再补集；或直接设置陷阱态。
\item 交集语言：先构造两个简单 DFA，再做乘积构造。
\end{itemize}


\subsection{1.3 正则运算和闭包}


正则语言对并、交、补、差、连接、Kleene star 封闭。


\begin{itemize}
\item 并/交/差/补：DFA 层面最直接，使用乘积或翻转接受态。
\item 连接/星号：NFA 层面最自然，使用 \(\epsilon\)-transition。
\item 反转 \(A^R\)：正则表达式结构归纳最直接：
\[
  (R_1\cup R_2)^R=R_1^R\cup R_2^R,\quad
  (R_1R_2)^R=R_2^RR_1^R,\quad
  (R^*)^R=(R^R)^*.
  \]

\end{itemize}


\subsection{1.4 Assignment 1 对齐}


官方答案要求会把图转为形式化描述。写 DFA 形式化描述时必须列：


\begin{itemize}
\item \(Q\)
\item \(\Sigma\)
\item \(\delta\) 表格
\item \(q_0\)
\item \(F\)
\end{itemize}


作业中的典型 DFA 语言：


\begin{itemize}
\item begins with \texttt{0} and ends with \texttt{1}
\item contains at least four \texttt{1}s
\item contains substring \texttt{0101}
\item length at least 4 and fourth symbol is \texttt{1}
\item starts with \texttt{1} and has odd length, or starts with \texttt{0} and has even length
\item contains at least two \texttt{0}s and at most one \texttt{1}
\item empty language
\item all strings except \(\epsilon\)
\end{itemize}


补集构造题：


\begin{itemize}
\item 不含 \texttt{ba}：先构造含 \texttt{ba} 的 DFA，再翻转接受态。
\item 既不含 \texttt{ab} 也不含 \texttt{ba}：等价于 \(a^*\cup b^*\)，也可先构造含 \texttt{ab} 或 \texttt{ba} 再补。
\item 不在 \(b^*a^*\) 中：等价于含 substring \texttt{ab}，注意方向不要写反。旧答案中题面与文字有编码混乱，复习时按答案源码的构造理解。
\end{itemize}


交集构造题：


\begin{itemize}
\item 至少两个 \texttt{a} 且至少三个 \texttt{b}：状态 \((i,j)\)，其中 \(i\in\{0,1,\ge2\}\)，\(j\in\{0,1,2,\ge3\}\)，接受 \((\ge2,\ge3)\)。
\item 偶数个 \texttt{a} 且一个或两个 \texttt{b}：状态记录 \texttt{a} parity 和 \texttt{b} count \(0,1,2,\ge3\)，接受 \texttt{(even,1)} 与 \texttt{(even,2)}。
\item 以 \texttt{a} 开头且至多两个 \texttt{b}：可以先判断首字符，再计数 \texttt{b}。
\end{itemize}


\bigskip\hrule\bigskip


\section{2. NFA、正则表达式、GNFA}


\subsection{2.1 NFA 定义}


NFA 仍是五元组，但转移函数为

\[
\delta:Q\times(\Sigma\cup\{\epsilon\})\to\mathcal P(Q).
\]

NFA 接受一个串，只要存在一条路径，路径标签（去掉 \(\epsilon\)）等于该串且终点在接受态。


\subsection{2.2 NFA 与 DFA 等价}


核心定理：若 NFA 识别 \(A\)，则 \(A\) 正则。证明用 subset construction。  

有 \(\epsilon\)-transition 时必须使用 \(\epsilon\)-closure。


\subsection{2.3 NFA 转 DFA：subset construction 完整写法}


考试或作业要求 “convert the following NFA to a DFA” 时，不能只写“由 subset construction 得”。答案要包含四件事：每个 DFA 状态代表的 NFA 状态集合、起始集合、转移表、接受集合。


给定

\[
N=(Q,\Sigma,\delta,q_0,F).
\]

若 \(N\) 有 \(\epsilon\)-transition，先定义

\[
E(R)=\{q\in Q\mid q\text{ can be reached from some }r\in R\text{ using only }\epsilon\text{-moves}\}.
\]

若没有 \(\epsilon\)-transition，可把 \(E(R)=R\)。


构造 DFA

\[
D=(Q_D,\Sigma,\delta_D,q_D,F_D)
\]

如下：


\begin{enumerate}
\item 起始态：
\[
   q_D=E(\{q_0\}).
   \]

\item 对每个已经发现的子集状态 \(S\subseteq Q\) 和每个输入符号 \(a\in\Sigma\)，计算
\[
   \operatorname{Move}(S,a)=\bigcup_{q\in S}\delta(q,a),
   \]

再令
\[
   \delta_D(S,a)=E(\operatorname{Move}(S,a)).
   \]

\item 若新集合还没出现，把它加入待处理队列。只列可达子集即可；若题目要求 complete DFA，则 \(\emptyset\) 也要作为死状态列出。
\item 接受态是所有包含旧接受态的子集：
\[
   F_D=\{S\in Q_D\mid S\cap F\neq\emptyset\}.
   \]

\end{enumerate}


正确性句子可以固定写：对任意输入串 \(w\)，DFA 读完 \(w\) 后所在的子集，恰好等于 NFA 从 \(q_0\) 读完 \(w\) 后可能处在的所有状态的 \(\epsilon\)-closure。因此该子集与 \(F\) 相交，当且仅当 NFA 有某条接受路径。


\subsection{2.4 NFA 转 DFA 示例 1：无 \(\epsilon\) 的 NFA}


构造一个识别“以 \texttt{01} 结尾”的 NFA：

\[
Q=\{q_0,q_1,q_2\},\quad \Sigma=\{0,1\},\quad q_0\text{ start},\quad F=\{q_2\}.
\]

转移为：

\[
\delta(q_0,0)=\{q_0,q_1\},\quad \delta(q_0,1)=\{q_0\},\quad
\delta(q_1,1)=\{q_2\},
\]

其它转移为空集。直觉是：\(q_0\) 任意循环，看到某个 \texttt{0} 时非确定性猜它是倒数第二个字符。


Subset construction 的可达集合如下。令

\[
A=\{q_0\},\quad B=\{q_0,q_1\},\quad C=\{q_0,q_2\}.
\]


\begin{longtable}{|p{0.22\textwidth}|p{0.22\textwidth}|p{0.22\textwidth}|p{0.22\textwidth}|}
\hline
\textbf{DFA state} & \textbf{on \texttt{0}} & \textbf{on \texttt{1}} & \textbf{accepting?} \\
\hline
\(A=\{q_0\}\) & \(B\) & \(A\) & no \\
\hline
\(B=\{q_0,q_1\}\) & \(B\) & \(C\) & no \\
\hline
\(C=\{q_0,q_2\}\) & \(B\) & \(A\) & yes \\
\hline
\end{longtable}


于是 DFA 为

\[
Q_D=\{A,B,C\},\quad q_D=A,\quad F_D=\{C\},
\]

转移表如上。注意 \(C\) 接受是因为 \(C\cap\{q_2\}\neq\emptyset\)，不是因为集合中所有 NFA 状态都接受。


\subsection{2.5 NFA 转 DFA 示例 2：带 \(\epsilon\)-closure}


设 NFA 识别语言 \(\{a,ab\}\)。状态

\[
Q=\{s,x,y,z,w,u\},\quad \Sigma=\{a,b\},\quad q_0=s,\quad F=\{y,u\}.
\]

转移为：

\[
\delta(s,\epsilon)=\{x,z\},\quad \delta(x,a)=\{y\},\quad
\delta(z,a)=\{w\},\quad \delta(w,b)=\{u\}.
\]

其它转移为空。先算起始闭包：

\[
E(\{s\})=\{s,x,z\}.
\]

令

\[
A=\{s,x,z\},\quad B=\{y,w\},\quad C=\{u\},\quad D=\emptyset.
\]


\begin{longtable}{|p{0.22\textwidth}|p{0.22\textwidth}|p{0.22\textwidth}|p{0.22\textwidth}|}
\hline
\textbf{DFA state} & \textbf{on \texttt{a}} & \textbf{on \texttt{b}} & \textbf{accepting?} \\
\hline
\(A=\{s,x,z\}\) & \(B\) & \(D\) & no \\
\hline
\(B=\{y,w\}\) & \(D\) & \(C\) & yes \\
\hline
\(C=\{u\}\) & \(D\) & \(D\) & yes \\
\hline
\(D=\emptyset\) & \(D\) & \(D\) & no \\
\hline
\end{longtable}


解释其中一行：

\[
\delta_D(A,a)=E(\delta(x,a)\cup\delta(z,a))=E(\{y,w\})=\{y,w\}=B.
\]

\(B\) 和 \(C\) 都接受，因为它们分别包含 \(y\) 或 \(u\)。这样得到的 DFA 精确接受 \texttt{a} 和 \texttt{ab}。


\subsection{2.6 正则表达式}


正则表达式递归定义：


\begin{itemize}
\item \(\emptyset,\epsilon,a\) 是正则表达式。
\item 若 \(R_1,R_2\) 是正则表达式，则 \(R_1\cup R_2\)、\(R_1R_2\)、\(R_1^*\) 是正则表达式。
\end{itemize}


Kleene theorem：

\[
\text{regular language}\iff \text{DFA/NFA}\iff \text{regular expression}.
\]


\subsection{2.7 Thompson 构造与状态消除}


\textbf{RE -> NFA}：对基础表达式建小 NFA，再递归处理 union/concat/star。


\textbf{DFA/NFA -> RE}：转 GNFA 后消状态。消去状态 \(q\) 时，对任意 \(i,j\) 更新边标签：

\[
R_{ij}\leftarrow R_{ij}\cup R_{iq}(R_{qq})^*R_{qj}.
\]

作业中必须写添加新 start、新 accept、逐步消状态，不能只给最终正则式。


\subsection{2.8 Assignment 2 对齐}


NFA 构造要点：


\begin{itemize}
\item ends with \texttt{00}：三态；起点可在任意处猜测倒数两个 \texttt{0} 的开始。
\item contains \texttt{0101}：五态；匹配进度 0 到 4。
\item even number of \texttt{0}s or exactly two \texttt{1}s：用新起点 \(\epsilon\) 分支到两个子 NFA。
\item \(\{0\}\)：两态，读一个 \texttt{0} 接受。
\item \(0^*1^*0^+\)：三态，注意最后的 \(0^+\) 至少一个 \texttt{0}。
\item \(1^*(001^+)^*\)：起点接受，读 \texttt{00} 后必须至少一个 \texttt{1} 回到起点。
\item \(\{\epsilon\}\)：单个起始接受态，无其它转移。
\item \(0^*\)：单个起始接受态，\texttt{0} 自环。
\end{itemize}


Closure construction 题必须按证明构造：


\begin{itemize}
\item Union：新起点 \(\epsilon\) 到两个起点。
\item Concatenation：第一个接受态 \(\epsilon\) 到第二个起点。
\item Star：新起点接受，并从旧接受态 \(\epsilon\) 回旧起点。
\item Single accept state：新增 \(q_a\)，旧接受态 \(\epsilon\) 到 \(q_a\)。
\end{itemize}


Subset construction 题要列子集状态表。答案源码里示例：


\begin{itemize}
\item 无 \(\epsilon\) 的 NFA：状态集形如 \(\emptyset,\{1\},\{2\},\{1,2\}\)。
\item 有 \(\epsilon\) 的 NFA：起始态是 \(\epsilon\)-closure，例如 \(\{1,2\}\)，接受态为包含旧接受态的子集。
\end{itemize}


Bonus：语言“不含一对相隔奇数个符号的 1”。答案思路是先构造补语言 NFA：猜第一个 \texttt{1}，用两个状态记录中间间隔奇偶，若在奇数间隔后看到第二个 \texttt{1} 则接受；再取补得到原语言 DFA。


\subsection{2.9 Assignment 3 对齐}


正则表达式速记：


\begin{itemize}
\item starts with \texttt{1}, ends with \texttt{0}：\(1\Sigma^*0\)
\item at least three \texttt{1}s：\(\Sigma^*1\Sigma^*1\Sigma^*1\Sigma^*\)
\item contains \texttt{0101}：\(\Sigma^*0101\Sigma^*\)
\item third symbol is \texttt{0}：\(\Sigma^2 0\Sigma^*\)
\item length at most 5：\((\Sigma\cup\epsilon)^5\)
\item strings of positive length：\(\Sigma^+\)
\item empty language：\(\emptyset\)
\end{itemize}


作业中的 state elimination 示例：


\begin{itemize}
\item 对两态 automaton，先加 \(s\) 和 \(f\)，若 1 到 2 为 \(b\)，2 自环 \(a\)，2 回 1 为 \(b\)，1 自环 \(a\)，最终答案形如
\[
  (a\cup ba^*b)^*ba^*.
  \]

\end{itemize}


注释语言 \(C\)：以 \texttt{/\#} 开始，以 \texttt{\#/} 结束，中间不能提前出现 \texttt{\#/}。答案给的正则式可简化记为：

\[
/\#(\#^*(a\cup b)\cup /)^*\#^+/
\]

关键状态含义：在 body 中看到 \texttt{\#} 后，若下一个是 \texttt{/} 就结束；若是 \texttt{a,b,\#} 则继续。


\bigskip\hrule\bigskip


\section{3. 非正则语言与 Pumping Lemma}


\subsection{3.1 正则 Pumping Lemma}


若 \(L\) 正则，则存在 \(p\)，任意 \(w\in L, |w|\ge p\)，可写为 \(w=xyz\)，满足：


\[
|xy|\le p,\quad |y|>0,\quad \forall i\ge0,\ xy^iz\in L.
\]


它只能证明“非正则”，不能证明“正则”。


\subsection{3.2 Assignment 4 标准证明}


**\(\{1^n2^n3^n\mid n\ge0\}\) 非正则**


取 \(w=1^p2^p3^p\)。由 \(|xy|\le p\)，\(y=1^\ell,\ell>0\)。泵 \(i=2\) 后得到 \(1^{p+\ell}2^p3^p\)，三段数量不相等，矛盾。


**\(\{www\mid w\in\{a,b\}^*\}\) 或答案中 \((a^pb^p)^3\) 类型**


取 \(w=(a^pb^p)^3\)。由于 \(|xy|\le p\)，\(y\) 位于第一段 \(a\) 中，泵后第一段变长，不能再分成三个完全相同的块，矛盾。


**\(\{a^{2^n}\mid n\ge0\}\) 非正则**


取 \(w=a^{2^p}\)，\(y=a^\ell\)，且 \(1\le\ell\le p\)。泵 \(i=2\) 后长度为 \(2^p+\ell\)，满足

\[
2^p<2^p+\ell\le2^p+p<2^{p+1}.
\]

长度夹在连续两个 2 的幂之间，矛盾。


**为什么不能用 \(0^p1^p\) 证明 \(0^*1^*\) 非正则**


要证明目标语言 \(0^*1^*\) 非正则，泵后必须离开 \(0^*1^*\)。但泵 \(0^p1^p\) 的前段 0 后仍是 \(0^*1^*\)。你证明的是不在 \(\{0^k1^k\}\)，不是不在 \(0^*1^*\)。


**\(\{0^m1^m0^m\mid m\ge0\}\) 非正则**


取 \(w=0^p1^p0^p\)。\(y\) 在第一段 0 内，泵 \(i=2\) 后前后 0 数量不同。


**\(\{0^m1^n\mid m\ne n\}\) 非正则**


用闭包：若 \(L\) 正则，则 \(\overline L\cap 0^*1^*\) 正则。但

\[
\overline L\cap 0^*1^*=\{0^n1^n\mid n\ge0\}
\]

非正则，矛盾。


**出现 \texttt{01} 次数等于 \texttt{10} 次数的语言是正则**


关键 claim：次数相等 iff \(w\in\{\epsilon,0,1\}\) 或首尾字符相同。  

正则表达式：

\[
\{\epsilon,0,1\}\cup0\Sigma^*0\cup1\Sigma^*1.
\]


**语言 \(C=\{1^k y\mid y\text{ 中 1 的个数}\le k\}\) 非正则**


取 \(w=1^p01^p\)，前缀 \(1^p\) 作 \(k=p\)。泵掉前面的 \(y=1^\ell\) 后，开头最多 \(p-\ell\) 个 1，但后缀仍有 \(p\) 个 1，违反“后缀 1 的个数 \(\le k\)”。


\bigskip\hrule\bigskip


\section{4. CFG、CFL、PDA}


\subsection{4.1 CFG 定义}


CFG 为

\[
G=(V,\Sigma,R,S)
\]

其中 \(V\) 是变量，\(\Sigma\) 是终结符，\(R\) 是产生式，\(S\) 是起始变量。  

若一个字符串有两棵不同 parse tree，则文法 ambiguous。


\subsection{4.2 常用 CFG}


Assignment 5 答案中的典型文法：


\begin{itemize}
\item 至少三个 \texttt{1}：
\[
  S\to R1R1R1R,\quad R\to0R\mid1R\mid\epsilon.
  \]

\item 以相同符号开始和结尾：
\[
  S\to1R1\mid0R0\mid0\mid1,\quad R\to0R\mid1R\mid\epsilon.
  \]

\item 奇数长度：
\[
  S\to0S0\mid0S1\mid1S0\mid1S1\mid0\mid1.
  \]

\item 奇数长度且中间符号为 0：
\[
  S\to0S0\mid0S1\mid1S0\mid1S1\mid0.
  \]

\item 回文：
\[
  S\to0S0\mid1S1\mid0\mid1\mid\epsilon.
  \]

\item 空语言：可用 \(S\to S\) 或无接受路径的 PDA。
\end{itemize}


\subsection{4.3 CNF}


Chomsky Normal Form 只允许：


\[
A\to BC,\quad A\to a
\]

如果语言含 \(\epsilon\)，允许新起始符号 \(S_0\to\epsilon\)。


转换步骤：


\begin{enumerate}
\item 加新起始变量 \(S_0\to S\)。
\item 消除 \(\epsilon\)-productions。
\item 消除 unit productions。
\item 长右部拆成二元变量，终结符替换成单独变量。
\end{enumerate}


Assignment 5 Q6 示例：


原文法：

\[
A\to BAB\mid B\mid\epsilon,\quad B\to11\mid\epsilon.
\]

加新起始符号，消 \(\epsilon\)，消 unit 后，引入 \(D\to1\), \(C\to AB\)，得到类似：

\[
A_0\to BC\mid BA\mid AB\mid BB\mid DD\mid\epsilon,
\quad A\to BC\mid BA\mid AB\mid BB\mid DD,
\quad C\to AB,\quad B\to DD,\quad D\to1.
\]


\subsection{4.4 PDA}


PDA 有输入、有限控制和栈。转移记为：

\[
a,X\to\gamma
\]

表示读入 \(a\)（可为 \(\epsilon\)），弹出 \(X\)（可为空，按课件记法），压入 \(\gamma\)。


核心等价：

\[
\text{CFL}\iff\text{某 PDA 识别}.
\]


Assignment 5 中正则语言对应的 PDA 可不使用栈；非正则 CFL 如 palindrome、奇数长度中间为 0 需要栈：


\begin{itemize}
\item palindrome：非确定性猜中点，前半段压栈，后半段逐字符匹配弹栈。
\item 奇数长度中间为 0：先压前半段计数，读中间 \texttt{0}，后半段每读一个符号弹一个计数符号。
\end{itemize}


\subsection{4.5 CFG 转 PDA：完整构造}


要证明每个 CFG 都有等价 PDA，答案必须写出 PDA 如何模拟最左推导。给定

\[
G=(V,\Sigma,R,S),
\]

构造 PDA \(P\)，使栈里保存“还需要匹配或展开的 sentential form”。复习时统一把栈顶写在最左边。


PDA 的状态可取

\[
Q=\{q_{\mathrm{start}},q,q_{\mathrm{acc}}\},
\]

栈字母表

\[
\Gamma=V\cup\Sigma\cup\{\$\}.
\]

核心转移如下：


\begin{enumerate}
\item 初始化：先把起始变量放到栈顶。
\[
   q_{\mathrm{start}}\xrightarrow{\epsilon,\$\to S\$}q.
   \]

\item 展开变量：对每条产生式
\[
   A\to X_1X_2\cdots X_k
   \]

加一条“用右部替换栈顶变量”的转移。若图或形式化定义只能一次压一个符号，则按
\[
   X_k,X_{k-1},\ldots,X_1
   \]

的顺序压栈，这样最后 \(X_1\) 在栈顶。简写为
\[
   q\xrightarrow{\epsilon,A\to X_1X_2\cdots X_k}q.
   \]

若 \(A\to\epsilon\)，则加
\[
   q\xrightarrow{\epsilon,A\to\epsilon}q.
   \]

\item 匹配终结符：对每个 \(a\in\Sigma\)，加
\[
   q\xrightarrow{a,a\to\epsilon}q.
   \]

意思是输入下一个符号和栈顶终结符相同，二者同时消掉。
\item 接受：输入读完且栈顶回到底符号时接受：
\[
   q\xrightarrow{\epsilon,\$\to\epsilon}q_{\mathrm{acc}}.
   \]

\end{enumerate}


正确性句子可以这样写：PDA 每次用 \(A\to u\) 替换栈顶变量，正好对应 CFG 最左推导中的一步；每次读入终结符 \(a\) 并弹出 \(a\)，正好确认当前推导出的最左终结符与输入一致。因此 PDA 恰好接受所有可由 \(S\) 推导出的字符串。


\subsection{4.6 CFG 转 PDA 示例：\(0^n1^n\)}


文法

\[
S\to 0S1\mid \epsilon
\]

生成

\[
\{0^n1^n\mid n\ge 0\}.
\]

按上面的构造，PDA 初始化时把 \(S\) 压到栈顶。展开规则为：

\[
\epsilon,S\to 0S1,\qquad \epsilon,S\to\epsilon.
\]

若只能逐个压栈，实现 \(S\to 0S1\) 时按 \(1,S,0\) 的顺序压，使得栈顶从左到右显示为 \(0S1\)。终结符匹配规则为：

\[
0,0\to\epsilon,\qquad 1,1\to\epsilon.
\]


输入 \texttt{0011} 的一次接受计算可以写成“栈内容变化”：

\[
S\$\Rightarrow 0S1\$\Rightarrow 00S11\$\Rightarrow 0011\$
\Rightarrow 011\$\Rightarrow 11\$\Rightarrow 1\$\Rightarrow \$.
\]

前三步是展开 \(S\)，后四步是读入并匹配输入。最后弹出 \(\$\) 进入接受态。


\subsection{4.7 PDA 转 CFG：完整构造}


要证明每个 PDA 都有等价 CFG，答案通常先把 PDA 规范化，再引入变量 \(A_{pq}\)。这是 Assignment 答案里最容易漏分的一类 construction proof。


先把 PDA 改成等价的规范形式：


\begin{itemize}
\item 只有一个接受态 \(q_{\mathrm{acc}}\)。
\item 接受前栈必须为空。
\item 每个转移只做一种栈操作：要么 push 一个栈符号，要么 pop 一个栈符号。若原转移替换多个符号，用中间状态拆开。
\end{itemize}


设规范化后的 PDA 状态集为 \(Q\)。为每一对状态 \(p,q\in Q\) 建一个 CFG 变量

\[
A_{pq}.
\]

含义必须写清楚：

\[
A_{pq}\text{ generates exactly the strings that take the PDA from }p
\text{ with empty stack to }q\text{ with empty stack}.
\]

起始变量是

\[
A_{q_0q_{\mathrm{acc}}}.
\]


产生式分三类：


\begin{enumerate}
\item 空计算：
\[
   A_{pp}\to\epsilon
   \]

对每个 \(p\in Q\) 都加入。
\item 串接计算：
\[
   A_{pq}\to A_{pr}A_{rq}
   \]

对所有 \(p,r,q\in Q\) 都加入。这表示从 \(p\) 到 \(q\) 的空栈计算可以在某个中间状态 \(r\) 分成两段。
\item push-pop 配对：
若 PDA 有
\[
   p\xrightarrow{a,\epsilon\to t}r
   \]

和
\[
   s\xrightarrow{b,t\to\epsilon}q,
   \]

其中 \(a,b\in\Sigma\cup\{\epsilon\}\)，则对所有这样的配对加入
\[
   A_{pq}\to aA_{rs}b.
   \]

解释：第一步 push 的 \(t\) 必须由最后一步 pop 掉，中间从 \(r\) 到 \(s\) 的计算不能动到这个 \(t\)，所以中间部分由 \(A_{rs}\) 生成。
\end{enumerate}


正确性句子可以固定写：任何从空栈到空栈的接受计算，要么是空计算，要么能在某个中间状态拆成两段，要么第一步 push 的栈符号与最后一次对应 pop 配对；三类产生式正好覆盖这三种情况。反过来，每条产生式都模拟 PDA 的合法计算片段，所以 CFG 不会生成额外字符串。


\subsection{4.8 PDA 转 CFG 示例：从 PDA 得到 \(S\to0S1\mid\epsilon\)}


构造一个规范 PDA 识别

\[
\{0^n1^n\mid n\ge0\}.
\]

状态为 \(p,r,s\)，从 \(p\) 开始，接受状态为 \(s\)，接受时栈为空。转移如下：

\[
p\xrightarrow{0,\epsilon\to X}p,\qquad
p\xrightarrow{\epsilon,\epsilon\to Y}r,\qquad
r\xrightarrow{\epsilon,Y\to\epsilon}s,\qquad
s\xrightarrow{1,X\to\epsilon}s.
\]

含义：在 \(p\) 中每读一个 \texttt{0} 就 push 一个 \(X\)；用 push \(Y\)、pop \(Y\) 的配对完成从 push 阶段到 pop 阶段的切换；在 \(s\) 中每读一个 \texttt{1} 就 pop 一个 \(X\)。


按构造，起始变量是 \(A_{ps}\)。由 \(Y\) 的 push-pop 配对：

\[
p\xrightarrow{\epsilon,\epsilon\to Y}r,\qquad
r\xrightarrow{\epsilon,Y\to\epsilon}s
\]

得到

\[
A_{ps}\to \epsilon A_{rr}\epsilon.
\]

又因为 \(A_{rr}\to\epsilon\)，所以有

\[
A_{ps}\to\epsilon.
\]

由 \(X\) 的 push-pop 配对：

\[
p\xrightarrow{0,\epsilon\to X}p,\qquad
s\xrightarrow{1,X\to\epsilon}s
\]

得到

\[
A_{ps}\to 0A_{ps}1.
\]

因此与起始变量有关的核心规则就是

\[
A_{ps}\to 0A_{ps}1\mid\epsilon,
\]

这正是生成 \(0^n1^n\) 的 CFG。考试中可以把 \(A_{ps}\) 重命名为 \(S\)，写成

\[
S\to0S1\mid\epsilon.
\]


\subsection{4.9 CFL 闭包与非闭包}


CFL 对并、连接、Kleene star 封闭。证明可用 CFG 或 PDA：


\begin{itemize}
\item 并：新起始变量 \(S\to S_1\mid S_2\)。
\item 连接：\(S\to S_1S_2\)。
\item Star：\(S\to S_1S\mid\epsilon\)。
\end{itemize}


每个正则语言都是 CFL：对正则表达式结构归纳，基础语言 \(\emptyset,\epsilon,\{a\}\) 都有 CFG，CFL 又对正则运算封闭。


CFL 不对交和补封闭：


\begin{itemize}
\item 取
\[
  A=\{a^mb^nc^n\mid m,n\ge0\},\quad
  B=\{a^nb^nc^m\mid m,n\ge0\}.
  \]

二者都是 CFL，但
\[
  A\cap B=\{a^nb^nc^n\mid n\ge0\}
  \]

不是 CFL。
\item 若 CFL 对补封闭，则由 De Morgan 与并封闭可推出对交封闭，矛盾。
\end{itemize}


CFL 与正则语言的交封闭：给 PDA \(P\) 和 DFA \(D\)，构造乘积 PDA，状态为 \((q_P,q_D)\)，栈行为照 PDA，读入符号时同时更新 DFA；接受态为 \(F_P\times F_D\)。


\subsection{4.10 CFL Pumping Lemma}


若 \(L\) 是 CFL，则存在 \(p\)，任意足够长 \(s\in L\) 可写为

\[
s=uvxyz
\]

满足

\[
|vxy|\le p,\quad |vy|\ge1,\quad \forall i\ge0,\ uv^ixy^iz\in L.
\]


Assignment 6 标准证明：


\begin{itemize}
\item \(\{0^n1^n0^n1^n\}\) 非 CFL：取 \(0^p1^p0^p1^p\)。因 \(|vxy|\le p\)，只影响至多两个相邻块。泵后至少一个块长度变，另有块保持 \(p\)，四段无法同为 \(n\)。
\item \(\{0^n\#0^{2n}\#0^{3n}\}\) 非 CFL：取 \(0^p\#0^{2p}\#0^{3p}\)。\(vxy\) 不可能跨两个 \texttt{\#}，只改变局部一段或相邻两段，泵后比例 \(1:2:3\) 被破坏。
\end{itemize}


\subsection{4.11 CNF 推导长度}


若 CFG 在 CNF 中，推导长度为 \(|w|\) 的非空串时恰好需要

\[
2|w|-1
\]

步：\(|w|-1\) 次 \(A\to BC\) 让变量叶子数从 1 增到 \(|w|\)，再 \(|w|\) 次 \(A\to a\) 生成终结符。


\bigskip\hrule\bigskip


\section{5. Turing Machines 与 Church-Turing}


\subsection{5.1 TM 定义}


标准 TM：

\[
M=(Q,\Sigma,\Gamma,\delta,q_0,q_{accept},q_{reject})
\]

其中 \(\Gamma\) 是 tape alphabet，\(\sqcup\in\Gamma\)，\(\Sigma\subseteq\Gamma-\{\sqcup\}\)，

\[
\delta:Q\times\Gamma\to Q\times\Gamma\times\{L,R\}.
\]


配置写作 \(uqv\)：带内容为 \(uv\)，当前状态 \(q\)，读写头在 \(v\) 的第一个字符上。


\subsection{5.2 Decider vs Recognizer}


\begin{itemize}
\item Decider：所有输入都 halt，接受 \(L\) 中输入，拒绝非 \(L\) 输入。
\item Recognizer：\(L\) 中输入会 accept；非 \(L\) 输入可以 reject 或 loop。
\item Decidable \(\Rightarrow\) Turing-recognizable。
\item \(A\) decidable iff \(A\) 和 \(\overline A\) 都 Turing-recognizable。
\end{itemize}


\subsection{5.3 TM 变体}


PPT 覆盖的等价模型：


\begin{itemize}
\item 多带 TM 等价于单带 TM。
\item NTM 等价于 DTM（可模拟所有计算分支）。
\item Enumerator 与 TM 识别器等价。
\item Church-Turing thesis：直观可算法计算的函数可由 TM 计算。
\end{itemize}


作业中重要变体：


**2-PDA 识别 \(\{a^nb^nc^n\}\)**  

两个栈即可模拟三段相等：读 \texttt{a} 时压 stack1；读 \texttt{b} 时压 stack2；读 \texttt{c} 时同时弹两个栈；输入结束且两栈空则接受。这说明 2-PDA 比 1-PDA 强。


\textbf{Left-reset TM 模拟普通左移}  

普通 TM 想从当前格左移时，left-reset machine 用标记法：


\begin{enumerate}
\item 标记当前格，reset 到最左端。
\item 从左端开始逐步找当前格前一格。
\item 用额外标记推进“候选前驱”。
\item 当候选右边是原当前格，擦标记并停在前驱格。
\end{enumerate}


答案强调：要写出这个算法过程，而不是只说“可以模拟”。


\textbf{只能右移/停留或输入只读的 TM}  

这类机器不能反复利用输入位置的无限记忆，识别能力退化为正则语言。证明思路：构造等价有限自动机，状态编码 TM 从当前位置向右移动后的有限控制状态。


\subsection{5.4 TM 设计算法题}


Assignment 7 典型题：


\begin{itemize}
\item 识别 \(\{a^nb^nc^n\}\)：循环扫描，找未标记 \texttt{a} 改为 \texttt{X}，向右找未标记 \texttt{b} 改为 \texttt{Y}，再找未标记 \texttt{c} 改为 \texttt{Z}，回到左端；最后确认没有未标记 \texttt{a,b,c} 且顺序合法。
\item \(\#0=\#1\)：反复找一个未标记 0 和一个未标记 1 配对标记；若最后只剩标记则接受。
\item \(\#0=2\#1\)：每找一个 \texttt{1}，匹配两个 \texttt{0}。
\item 补语言：若已有 decider，可翻转 accept/reject；若只是 recognizer，不能直接补。
\item 二进制减一：移到最右，从右向左把连续 0 改成 1，遇到第一个 1 改成 0 后停。
\item 二进制减法 \(x-y\)：重复对 \(x\) 和 \(y\) 执行减一，直到 \(y=0\)。
\item 整数除法 \(\lfloor x/y\rfloor\)：多带 TM 上重复减 \(y\)，计数 quotient。
\end{itemize}


\bigskip\hrule\bigskip


\section{6. Decidability 与 Undecidability}


\subsection{6.1 可判定语言}


经典可判定问题：


\begin{itemize}
\item \(A_{DFA}=\{\langle M,w\rangle\mid M\text{ is a DFA and accepts }w\}\)
\item \(A_{NFA}\)
\item \(A_{REX}\)
\item \(E_{DFA}=\{\langle M\rangle\mid L(M)=\emptyset\}\)
\item \(EQ_{DFA}\)
\item \(A_{CFG}\)
\item \(E_{CFG}\)
\end{itemize}


典型算法：


\begin{itemize}
\item \(A_{DFA}\)：直接模拟 DFA。
\item \(A_{NFA}\)：转 DFA 或图搜索所有 NFA 分支。
\item \(A_{REX}\)：RE -> NFA -> DFA，再判定。
\item \(E_{DFA}\)：从 start 做 BFS，看是否能到接受态。
\item \(EQ_{DFA}\)：构造对称差 DFA，判空。
\item \(A_{CFG}\)：转 CNF 后 CYK/dynamic programming。
\item \(E_{CFG}\)：标记能生成终结串的变量，看 start 是否被标记。
\end{itemize}


\subsection{6.2 Assignment 8 对齐}


\textbf{DFA 与 REX 等价可判定}


对输入 \(\langle M,R\rangle\)：


\begin{enumerate}
\item 将 \(R\) 转为 NFA。
\item 将 NFA 转为 DFA \(M'\)。
\item 运行 \(EQ_{DFA}\) 于 \(\langle M,M'\rangle\)。
\end{enumerate}


**\(ALL_{DFA}\) 可判定**


构造识别 \(\overline{L(A)}\) 的 DFA \(M\)，运行 \(E_{DFA}\)。若补语言为空，则 \(L(A)=\Sigma^*\)。


**\(A_{\epsilon CFG}\) 可判定**


运行 \(A_{CFG}\) 于 \(\langle G,\epsilon\rangle\)。或转 CNF 后检查是否有 \(S\to\epsilon\)。


**\(\overline{E_{TM}}\) Turing-recognizable**


对输入 \(\langle M\rangle\)，枚举所有字符串 \(s_1,s_2,\ldots\)，dovetail 模拟：第 \(i\) 轮让前 \(i\) 个字符串各跑 \(i\) 步。一旦某个输入被接受，就接受。若 \(L(M)=\emptyset\)，机器可能永远不接受。


**\(EQ_{DFA}\) 有限测试上界**


若 \(A,B\) 有 \(n_1,n_2\) 个状态，则对称差 DFA 有 \(n_1n_2\) 个状态。若存在区分串，则存在长度小于 \(n_1n_2\) 的区分串。需测试数量：

\[
\sum_{i=0}^{n_1n_2-1}|\Sigma|^i.
\]


\subsection{6.3 Diagonalization 与 \(A_{TM}\)}


\[
A_{TM}=\{\langle M,w\rangle\mid M\text{ accepts }w\}
\]

是 Turing-recognizable 但 undecidable。


Recognizer：通用 TM 模拟 \(M(w)\)，若 \(M\) accept 则 accept，若 reject 则 reject，若 loop 则 loop。


不可判定证明：


\begin{enumerate}
\item 假设 \(H\) decides \(A_{TM}\)。
\item 构造 \(D\)：输入 \(\langle M\rangle\)，运行 \(H(\langle M,\langle M\rangle\rangle)\)。若 \(H\) accept，则 \(D\) reject；若 \(H\) reject，则 \(D\) accept。
\item 运行 \(D(\langle D\rangle)\) 得矛盾。
\end{enumerate}


\(\overline{A_{TM}}\) 不是 Turing-recognizable。否则 \(A_{TM}\) 和其补都 recognizable，可并行运行得到 decider，矛盾。


\subsection{6.4 Reducibility}


映射归约：

\[
A\le_m B
\]

若存在可计算函数 \(f\)，使得

\[
w\in A\iff f(w)\in B.
\]


性质：


\begin{itemize}
\item 若 \(A\le_m B\) 且 \(B\) decidable，则 \(A\) decidable。
\item 若 \(A\le_m B\) 且 \(A\) undecidable，则 \(B\) undecidable。
\item 若 \(A\le_m B\) 且 \(B\) T-recognizable，则 \(A\) T-recognizable。
\item 若 \(A\le_m B\) 且 \(A\) T-unrecognizable，则 \(B\) T-unrecognizable。
\end{itemize}


\subsection{6.5 Rice Theorem}


任何关于 \(L(M)\) 的非平凡性质都是 undecidable。  

非平凡：有些 TM 的语言满足该性质，有些不满足。  

注意：Rice 只适用于“语言性质”，不适用于“TM 语法性质”如状态数是否大于 481。


Assignment 10 中用 Rice 证明：


\begin{itemize}
\item \(L(M)\) infinite：非平凡语言性质，undecidable。
\item \(1011\in L(M)\)：非平凡语言性质，undecidable。
\item \(L(M)=\Sigma^*\)：非平凡语言性质，undecidable。
\end{itemize}


\subsection{6.6 Assignment 9-10 对齐}


\textbf{不可数性}


\(\{0,1,2\}^{\mathbb N}\) 不可数：假设列成 \(f(0),f(1),\ldots\)，构造新序列 \(s\)，第 \(i\) 位取不同于 \(f(i)\) 第 \(i\) 位的符号。则 \(s\) 与每一行在对角位不同。


\textbf{三元组集合可数}


按 \(i+j+k=n\) 分层，每层有限：

\[
\binom{n+2}{2}
\]

先列和小的，再列同层内部，得到枚举。


\textbf{“TM 是否会在空白带上写非空符号”可判定}


如果在写非空符号前，带一直全 blank，行为只由状态决定。模拟 \(|Q|+1\) 步：若没写且状态重复，则之后永远重复，不会写。


\textbf{“TM 是否至少 481 个状态”可判定}


这是语法性质，直接解析编码并计数状态。


\textbf{DFA 是否不接受任何偶数个 1 的串}


构造 DFA \(A\) 接受偶数个 1 的串，构造 \(M\cap A\)，运行 \(E_{DFA}\) 判空。


**\(\overline{EQ_{CFG}}\) Turing-recognizable**


枚举所有字符串 \(w\)，运行 \(A_{CFG}\) 判断 \(w\) 是否由 \(G_1,G_2\) 生成；若结果不同则接受。若两个 CFG 等价，则永不接受。因此 \(EQ_{CFG}\) 是 co-Turing-recognizable。


**若 \(A\) T-recognizable 且 \(A\le_m\overline A\)，则 \(A\) decidable**


由归约 \(f\) 得：

\[
w\in\overline A \iff f(w)\in A.
\]

因为 \(A\) recognizable，先计算 \(f(w)\)，再运行 \(A\) 的 recognizer，即可 recognize \(\overline A\)。所以 \(A\) 与 \(\overline A\) 都 recognizable，故 \(A\) decidable。


\textbf{判定 halts on empty input 不可判定}


从 \(HALT_{TM}\) 归约。给 \(\langle M,w\rangle\)，构造 \(M'\)：在空输入上模拟 \(M(w)\)，若 \(M(w)\) halt 则 halt，否则 loop。若能判定 \(M'\) 是否在空输入 halt，就能判定 \(HALT_{TM}\)。


\bigskip\hrule\bigskip


\section{7. Time Complexity 与 P}


\subsection{7.1 时间复杂度}


TM 在时间 \(t(n)\) 内运行：对所有长度 \(n\) 输入，最多 \(t(n)\) 步内 halt。


\[
TIME(t(n))=\{B\mid B\text{ 可由确定性单带 TM 在 }O(t(n))\text{ 时间判定}\}.
\]


\[
P=\bigcup_k TIME(n^k)
\]

即多项式时间可判定语言。


课件强调：


\begin{itemize}
\item 最坏情况复杂度是对长度 \(n\) 的所有输入取上界。
\item reasonable encoding 很重要；图的 adjacency matrix/list 是合理编码，数字用 unary 通常不合理。
\item 多带 TM 与单带 TM 在多项式意义下等价：多带 \(t(n)\) 可由单带 \(O(t(n)^2)\) 模拟。
\end{itemize}


\subsection{7.2 P 中的典型问题}


**PATH \(\in P\)**  

给有向图 \(G,s,t\)，从 \(s\) 做 BFS/marking，若标记到 \(t\) 则接受。


**RELPRIME \(\in P\)**  

用 Euclidean algorithm：重复 \(x\leftarrow x\bmod y\)，交换 \(x,y\)，直到 \(y=0\)。若 gcd 为 1 则接受。


**CFL \(\in P\)**  

用 CYK/dynamic programming 而不是暴力枚举推导。转 CNF 后，对每个 substring 存能生成它的变量集合，总复杂度多项式，典型为 \(O(n^3)\) 级别乘文法大小因子。


\textbf{HAMPATH 与 PATH 的区别}  

PATH 只问是否有任意路径，BFS 可解；HAMPATH 要经过每个节点恰好一次，暴力路径数可达 \(m!\)，是否在 P 是开放/与 NP 完全性相关。


\subsection{7.3 P 闭包与 Assignment 10}


P 对并、连接、补封闭。


\begin{itemize}
\item Union：顺序运行两个多项式 decider，任一接受则接受。
\item Concatenation：枚举 \(n+1\) 个 split \(w=uv\)，运行两个 decider，某个 split 都接受则接受。
\item Complement：运行 decider 并翻转结果。
\end{itemize}


**CONNECTED \(\in P\)**  

无向图从任意顶点 BFS，若所有顶点都被标记则接受。


**ALL\_DFA \(\in P\)**  

构造补 DFA，对补 DFA 从起点 BFS。若能到接受态，则原 DFA 不是 all；否则原 DFA 接受 \(\Sigma^*\)。


**EQ\_DFA \(\in P\)**  

构造对称差 DFA，再用 \(E_{DFA}\) 的 BFS 判空。


\bigskip\hrule\bigskip


\section{8. NP、coNP、Polynomial Reducibility}


\subsection{8.1 NP 定义}


\[
NP=\bigcup_k NTIME(n^k)
\]

等价定义：存在多项式时间 verifier \(V\) 和多项式长度 certificate \(c\)，使得


\[
w\in L\iff \exists c,\ |c|\le |w|^k,\ V(w,c)\text{ accepts}.
\]


PPT 中重要定理：语言在 NP 中 iff 有 polynomial-time verifier。  

NTM -> verifier：certificate 描述接受分支。  

verifier -> NTM：非确定性猜 certificate，再运行 verifier。


\subsection{8.2 常见 NP 证明}


Assignment 11 要求证明这些语言在 NP。标准写法是给 certificate 和 verifier：


\begin{itemize}
\item 3SAT：certificate 是 truth assignment；检查每个 clause 是否有 true literal。
\item VERTEX-COVER：certificate 是 \(k\) 个顶点；检查每条边至少一个端点在集合中。
\item TSP：certificate 是 tour；检查每个城市恰好一次且总权重 \(\le k\)。
\item MAX-CUT：certificate 是顶点划分 \(S,V-S\)；数 crossing edges 是否 \(\ge k\)。
\item 3-COLORING：certificate 是每个顶点颜色；检查相邻顶点颜色不同。
\item CLIQUE：certificate 是 \(k\) 个顶点；检查任意两点之间有边。
\item SUBSET-SUM：certificate 是子集；检查和是否为目标 \(t\)。
\item COMPOSITES：certificate 是非平凡因子 \(y\)；检查 \(1<y<x\) 且 \(y\mid x\)。
\end{itemize}


\subsection{8.3 NP 闭包}


NP 对 union、concatenation、star 封闭。


\begin{itemize}
\item Union：certificate 额外包含选择 bit，说明用 \(V_1\) 还是 \(V_2\)。
\item Concatenation：certificate 包含 split 位置 \(i\) 以及两个子证书 \(c_1,c_2\)。
\item Star：certificate 包含分割位置列表和每段证书；段数至多 \(|w|\)，总长度仍多项式。
\end{itemize}


\subsection{8.4 P、NP、coNP 判断题}


标准结论：


\begin{itemize}
\item \(P\subseteq NP\)：True，忽略 certificate，直接运行 P decider。
\item \(P\) 对补封闭：True。
\item \(P\subseteq NP\cap coNP\)：True。
\item 若 \(P=NP\)，则 \(NP=coNP\)：True。
\item \(\overline{SAT}\in P\)：严格说是开放问题；若它在 P，则 \(P=NP\)。作业官方 revised answer 更严谨地说明 unknown；旧简答里写 False 是基于通常相信 \(P\ne NP\)，考试应按题目/答案口径说明“unknown unless assuming \(P\ne NP\)”。
\end{itemize}


\subsection{8.5 2-COLOR 与 2-SAT}


**2-COLOR \(\in P\cap NP\)**  

NP：certificate 是 2-coloring。  

P：BFS/DFS 检查二分图。未染色点染 0，邻居染 1，若遇到同色边则 reject。


**2-SAT \(\in P\cap NP\)**  

NP：certificate 是 assignment。  

P：构造 implication graph：

\[
(x\vee y)\equiv(\neg x\to y)\wedge(\neg y\to x).
\]

若某变量 \(x\) 与 \(\neg x\) 在同一个 SCC，则不可满足；否则可满足。


\bigskip\hrule\bigskip


\section{9. NP-Completeness}


\subsection{9.1 定义与证明套路}


语言 \(B\) 是 NP-complete iff：


\begin{enumerate}
\item \(B\in NP\)
\item 对所有 \(A\in NP\)，\(A\le_p B\)
\end{enumerate}


实际证明 \(C\) NP-complete：


\begin{enumerate}
\item 证明 \(C\in NP\)。
\item 从一个已知 NP-complete 问题 \(B\) 归约到 \(C\)，如 \(3SAT\le_p C\)。
\item 证明构造多项式时间。
\item 证明 iff。
\end{enumerate}


\subsection{9.2 3SAT -> CLIQUE}


给 3CNF：

\[
\phi=C_1\wedge\cdots\wedge C_m.
\]

构造图 \(G\)：


\begin{enumerate}
\item 对每个 clause 中每个 literal occurrence 建一个顶点。
\item 不同 clause 的两个顶点之间连边，当且仅当两个 literal 不矛盾。
\item 设置 \(k=m\)。Assignment 12 Q1 中 \(m=3\)，所以 \(k=3\)。
\end{enumerate}


正确性：


\begin{itemize}
\item 若 \(\phi\) satisfiable，从每个 clause 选一个 true literal。这些 literal 互不矛盾，所以对应顶点两两相连，形成 \(m\)-clique。
\item 若 \(G\) 有 \(m\)-clique，因为同 clause 内没有边，clique 必定每个 clause 选一个顶点；两两相连说明 literal 互不矛盾，可扩展为满足赋值。
\end{itemize}


\subsection{9.3 2SAT -> CLIQUE}


同样构造，每个 2-clause 建两个顶点，连接不同 clause 且不矛盾的 literal，令 \(k=m\)。  

能说明 2SAT 实例可多项式变成 CLIQUE 实例；但由于 2SAT 在 P，这个归约不能证明 CLIQUE NP-hard。要证明 CLIQUE NP-hard 必须从 NP-complete 问题如 3SAT 归约。


\subsection{9.4 Double-SAT NP-complete}


Double-SAT：

\[
\{\phi\mid \phi\text{ has at least two satisfying assignments}\}.
\]


在 NP：certificate 是两个不同 assignment，验证二者不同且都满足 \(\phi\)。


从 3SAT 归约：给 \(\phi\)，引入新变量 \(w\)，构造

\[
\phi'=\phi\wedge(w\vee\overline w).
\]

若 \(\phi\) satisfiable，则任意满足赋值可扩展为 \(w=true\) 和 \(w=false\) 两个满足赋值。若 \(\phi'\) 有至少两个满足赋值，则特别地 \(\phi'\) satisfiable，从而 \(\phi\) satisfiable。


\subsection{9.5 3SAT -> HAMPATH}


课件与作业采用 Sipser diamond variable gadget：


\begin{itemize}
\item 每个变量一个 gadget。
\item 从左到右遍历表示 true，从右到左遍历表示 false。
\item 每个 clause 一个 clause vertex，通过 wire 连到对应 literal 在 gadget 中的位置。
\item 若某 literal 使 clause satisfied，Hamiltonian path 可绕入该 clause vertex 并返回。
\end{itemize}


证明结构：


\begin{itemize}
\item 满足赋值 -> 按变量真值选择每个 gadget 的遍历方向，对每个 clause 通过一个 true literal 的 wire 插入 clause vertex，得到 Hamiltonian path。
\item Hamiltonian path -> 每个变量 gadget 的方向给出赋值；每个 clause vertex 被访问必须通过某个 literal wire，该 literal 在对应方向下为 true，所以每个 clause satisfied。
\end{itemize}


\subsection{9.6 Directed HAMPATH -> Undirected HAMPATH}


Assignment 13 Q4：


给 directed \(G=(V,E)\)，构造 undirected \(G'\)。对每个顶点 \(u\) 建：

\[
u^{in},u^{mid},u^{out}
\]

并加内部边：

\[
\{u^{in},u^{mid}\},\quad \{u^{mid},u^{out}\}.
\]

对每条有向边 \((u,v)\)，加无向边：

\[
\{u^{out},v^{in}\}.
\]

起点终点：

\[
s'=s^{in},\quad t'=t^{out}.
\]


关键证明：\(u^{mid}\) 度只连 \(u^{in},u^{out}\)，所以 Hamiltonian path 必须连续穿过一个 gadget。路径从 \(s^{in}\) 开始会强制所有 gadget 都以 \(in\to mid\to out\) 方向穿过；跨 gadget 的边只能对应原图有向边。


\subsection{9.7 3SAT -> VERTEX-COVER}


Assignment 13 Q5：


给 3CNF \(\phi\)，变量数 \(v\)，clause 数 \(c\)。


构造：


\begin{enumerate}
\item 每个变量 \(x_i\)：建两个顶点 \(x_i,\overline{x_i}\)，加边 \(\{x_i,\overline{x_i}\}\)。
\item 每个 clause \(C_j=(\ell_{j,1}\vee\ell_{j,2}\vee\ell_{j,3})\)：建三角形三个 clause 顶点。
\item 每个 clause literal 顶点连到变量 gadget 中同名 literal 顶点。
\item 设置
\[
   k=v+2c.
   \]

\end{enumerate}


正确性：


\begin{itemize}
\item 满足赋值 -> 每个变量 gadget 选 true literal 顶点；每个 clause 三角形中留下一个 true literal 顶点不选，选另两个。变量边、三角形边、连接边都被覆盖，大小 \(v+2c\)。
\item Vertex cover 大小 \(\le v+2c\) -> 每个变量边至少选一个，共至少 \(v\)；每个三角形至少选两个，共至少 \(2c\)。因此必须恰好选这些数量。每个 clause 三角形有一个未选顶点，其连接边必须由变量 gadget 中同名顶点覆盖，令该 literal 为 true，即每个 clause 有 true literal。
\end{itemize}


\subsection{9.8 3COLOR NP-complete}


Assignment 12 Q6：


在 NP：certificate 是 coloring，逐边检查端点颜色不同。


从 3SAT 归约：


\begin{enumerate}
\item 建 palette 三角形 \(T,F,N\)，强制三种颜色。
\item 每个变量建 \(x_i,\overline{x_i}\)，二者互连，且都连到 \(N\)。因此它们不能取 Neutral，且必须一真一假。
\item 每个 clause \((a\vee b\vee c)\) 用两个 OR gadget：先算 \(u=a\vee b\)，再算 \(o=u\vee c\)。
\item gadget 的中间输出和最终输出连到 \(N\)，使其只能 True/False。
\item 最终输出 \(o\) 还连到 \(F\)，强制 \(o\) 为 True。
\end{enumerate}


正确性：clause 三个 literal 全 false 时 OR output 被迫 false，与连到 \(F\) 冲突；至少一个 true 时 gadget 可合法 3-color。


\bigskip\hrule\bigskip


\section{10. Cook-Levin 与 3SAT}


\subsection{10.1 Cook-Levin Theorem}


\[
SAT\text{ is NP-complete}.
\]


证明结构：


\begin{enumerate}
\item \(SAT\in NP\)：certificate 是 truth assignment。
\item 对任意 \(A\in NP\)，取多项式时间 NTM \(M\) 决定 \(A\)。给输入 \(w\)，构造公式 \(\phi_{M,w}\)，使
\[
   w\in A\iff \phi_{M,w}\text{ satisfiable}.
   \]

\end{enumerate}


\subsection{10.2 Tableau}


若 \(M\) 在 \(n^k\) 时间内运行，构造 \(n^k\times n^k\) tableau，每行是一条 configuration。变量：

\[
x_{i,j,s}
\]

表示 tableau 的第 \(i\) 行第 \(j\) 格内容是符号/状态 \(s\)。


公式：

\[
\phi_{M,w}=\phi_{cell}\wedge\phi_{start}\wedge\phi_{move}\wedge\phi_{accept}.
\]


\begin{itemize}
\item \(\phi_{cell}\)：每个 cell 恰好一个符号。
\item \(\phi_{start}\)：第一行是 \(q_0w\sqcup\cdots\)。
\item \(\phi_{accept}\)：某处出现 \(q_{accept}\)。
\item \(\phi_{move}\)：每个 \(2\times3\) window 合法，保证相邻配置符合转移函数。
\end{itemize}


\subsection{10.3 为什么用 \(2\times3\) window}


Assignment 13 Q3 反例：


设

\[
\delta(q,0)=(r,0,L),\quad \delta(q,1)=(s,1,L),\quad r\ne s.
\]

考虑 window：

\[
\begin{array}{ccc}
0&q&1\\
r&0&1
\end{array}
\]

它非法，因为上方状态 \(q\) 实际扫描的是右侧 \texttt{1}，应产生 \(s\) 而不是 \(r\)。但它的两个 \(2\times2\) 子窗口都可分别出现在某个合法 \(2\times3\) window 中。因此只查 \(2\times2\) 不够。


\subsection{10.4 Assignment 13 Q1 tableau}


TM 识别 odd strings，对 \(w=111\) 的计算：

\[
q_0 111\bot \vdash 1q_1 11\bot \vdash 11q_0 1\bot
\vdash 111q_1\bot \vdash 111\bot q_{acc}.
\]

对应 \(5\times5\) tableau：


\begin{longtable}{|p{0.176\textwidth}|p{0.176\textwidth}|p{0.176\textwidth}|p{0.176\textwidth}|p{0.176\textwidth}|}
\hline
\textbf{1} & \textbf{2} & \textbf{3} & \textbf{4} & \textbf{5} \\
\hline
\(q_0\) & 1 & 1 & 1 & \(\bot\) \\
\hline
1 & \(q_1\) & 1 & 1 & \(\bot\) \\
\hline
1 & 1 & \(q_0\) & 1 & \(\bot\) \\
\hline
1 & 1 & 1 & \(q_1\) & \(\bot\) \\
\hline
1 & 1 & 1 & \(\bot\) & \(q_{acc}\) \\
\hline
\end{longtable}


\[
\phi_{start}=x_{1,1,q_0}\land x_{1,2,1}\land x_{1,3,1}\land x_{1,4,1}\land x_{1,5,\bot}.
\]

\[
\phi_{accept}=\bigvee_{i=1}^{5}\bigvee_{j=1}^{5}x_{i,j,q_{acc}}.
\]


\subsection{10.5 SAT -> 3SAT}


3SAT 是 NP-complete。长 clause 转 3CNF：


\begin{itemize}
\item 1 literal：\(x\) 可写成 \((x\vee x\vee x)\)。
\item 2 literals：\((x\vee y)\) 可写成 \((x\vee y\vee x)\) 或用 padding。
\item 3 literals：保持。
\item \(l>3\) literals：
\[
  (x_1\vee\cdots\vee x_l)
  \]

替换为
\[
  (x_1\vee x_2\vee z_1)\wedge(\overline z_1\vee x_3\vee z_2)\wedge\cdots
  \wedge(\overline z_{l-3}\vee x_{l-1}\vee x_l).
  \]

\end{itemize}

注意课件说原公式与新公式不必逻辑等价，只需 satisfiability 等价。


\bigskip\hrule\bigskip


\section{11. Space, PSPACE, NP-hardness, 近似}


\subsection{11.1 Space complexity}


TM 在 space \(f(n)\) 内运行：对长度 \(n\) 输入，最多使用 \(f(n)\) 个 tape cells。


\[
SPACE(f(n))=\{B\mid B\text{ 被某 deterministic TM 用 }O(f(n))\text{ space 判定}\}
\]

\[
NSPACE(f(n))=\{B\mid B\text{ 被某 nondeterministic TM 用 }O(f(n))\text{ space 判定}\}.
\]


\[
PSPACE=\bigcup_k SPACE(n^k),\quad NPSPACE=\bigcup_k NSPACE(n^k).
\]


Savitch theorem：

\[
NSPACE(f(n))\subseteq SPACE(f(n)^2),\quad f(n)\ge\log n.
\]

推论：

\[
PSPACE=NPSPACE.
\]


课件关系：

\[
P\subseteq NP\subseteq PSPACE\subseteq EXPTIME
\]

且 \(P\ne EXPTIME\)，所以这些包含中至少一个是真包含，但不知道是哪一个。


\subsection{11.2 NP-hardness}


NP-hard 不要求问题在 NP 中，也可用于 search/optimization 问题。  

定义：对所有 \(A\in NP\)，\(A\le_p B\)。若某 NP-hard 问题有多项式算法，则 \(P=NP\)。


NP-complete = NP-hard + in NP。


\subsection{11.3 近似与随机算法}


\textbf{Vertex Cover 2-approximation}


算法：当图中还有边时，任选一条边 \((u,v)\)，把 \(u,v\) 都加入 cover，并删除所有 incident edges。


证明：


\begin{enumerate}
\item 输出 \(X\) 是 vertex cover，因为每条边在被处理或删除时都被选中端点覆盖。
\item 选中的边彼此不共享端点，任意最优 cover 至少要为每条这样的边选一个端点。
\item 算法每条选中边选两个端点，所以 \(|X|\le2OPT\)。
\end{enumerate}


**Randomized MAX-3SAT \(7/8\) expectation**


每个变量独立以 \(1/2\) 取 true。一个 3-literal clause 不满足概率为

\[
(1/2)^3=1/8
\]

满足概率 \(7/8\)。令 \(X_i\) 是第 \(i\) 个 clause 是否满足的 indicator，则

\[
\mathbb E[X_i]=7/8,\quad
\mathbb E[X]=\sum_i\mathbb E[X_i]=7m/8.
\]


\bigskip\hrule\bigskip


\section{12. 快速总表}


\subsection{12.1 模型能力}


\begin{longtable}{|p{0.293\textwidth}|p{0.293\textwidth}|p{0.293\textwidth}|}
\hline
\textbf{模型} & \textbf{等价描述} & \textbf{语言类} \\
\hline
DFA & NFA, regex & Regular \\
\hline
PDA & CFG & CFL \\
\hline
TM decider & 总停机算法 & Decidable \\
\hline
TM recognizer & 接受时停机 & Turing-recognizable \\
\hline
Poly-time DTM & efficient decider & P \\
\hline
Poly-time NTM / verifier & short certificate & NP \\
\hline
\end{longtable}


\subsection{12.2 常见“可判定/不可判定/可识别”}


\begin{longtable}{|p{0.44\textwidth}|p{0.44\textwidth}|}
\hline
\textbf{问题} & \textbf{结论} \\
\hline
\(A_{DFA},A_{NFA},A_{REX}\) & decidable \\
\hline
\(E_{DFA},EQ_{DFA}\) & decidable \\
\hline
\(A_{CFG},E_{CFG}\) & decidable \\
\hline
\(A_{TM}\) & recognizable, undecidable \\
\hline
\(\overline{A_{TM}}\) & not recognizable \\
\hline
\(HALT_{TM}\) & undecidable \\
\hline
\(E_{TM}\) & undecidable, not recognizable \\
\hline
\(\overline{E_{TM}}\) & recognizable \\
\hline
\(EQ_{TM}\) & undecidable, not recognizable \\
\hline
\(\overline{EQ_{CFG}}\) & recognizable \\
\hline
\end{longtable}


\subsection{12.3 常见 NP-complete 链}


\[
SAT \le_p 3SAT \le_p CLIQUE
\]

\[
3SAT \le_p HAMPATH
\]

\[
HAMPATH \le_p UHAMPATH
\]

\[
3SAT \le_p VERTEX\text{-}COVER
\]

\[
3SAT \le_p 3COLOR
\]


\subsection{12.4 作业对应复习索引}


\begin{longtable}{|p{0.44\textwidth}|p{0.44\textwidth}|}
\hline
\textbf{作业} & \textbf{必会内容} \\
\hline
Assignment 1 & DFA 形式化、补集、乘积构造 \\
\hline
Assignment 2 & NFA 图、union/concat/star construction、single accept、subset construction \\
\hline
Assignment 3 & regex、Thompson、GNFA state elimination、reverse 正则性 \\
\hline
Assignment 4 & 正则 pumping lemma、闭包反证 \\
\hline
Assignment 5 & CFG、parse tree、PDA、CNF \\
\hline
Assignment 6 & CFG-PDA 等价、CFL closure、CFL pumping、CFL 非闭包、CFL $\cap$ regular \\
\hline
Assignment 7 & 2-PDA、TM 配置、TM 设计、left-reset、read-only/right-only 正则性、decidable/TR closure \\
\hline
Assignment 8 & DFA/RE/CFG 判定问题、\(\overline{E_{TM}}\) recognizer、\(EQ_{DFA}\) 长度上界 \\
\hline
Assignment 9 & 对角线、可数性、TM 语法/行为判定、DFA 与偶数个 1 \\
\hline
Assignment 10 & \(EQ_{CFG}\) co-recognizable、mapping reduction、Rice、CYK、P 闭包、CONNECTED/ALL\_DFA/EQ\_DFA in P \\
\hline
Assignment 11 & NP verifier、NP closure、2COLOR、2SAT、P/NP/coNP 判断 \\
\hline
Assignment 12 & 3SAT->CLIQUE、2SAT->CLIQUE、Double-SAT、3SAT->HAMPATH、3COLOR \\
\hline
Assignment 13 & Cook-Levin tableau、legal window、UHAMPATH、VERTEX-COVER \\
\hline
\end{longtable}


\bigskip\hrule\bigskip


\section{13. 最后复习提醒}


\begin{enumerate}
\item 自动机题：先写“状态记录什么”，再写状态/转移/接受态。
\item 构造证明题：必须给完整 tuple 或完整规则，不要只说“按闭包性”。
\item Pumping 题：选择的串必须在语言中，且分解必须任意。
\item 归约题：方向绝对不能反。从已知难问题归约到目标问题。
\item NP-complete 题：先 in NP，再 NP-hard。
\item Rice 题：先确认是 \(L(M)\) 的非平凡语义性质，不是机器编码的语法性质。
\item P/NP 题：注意“可验证”与“可求解”不同；HAMPATH 在 NP，但不知道是否在 P。
\end{enumerate}

\end{document}
